%% Derived from templates/daily.tex + weekly.tex on 2026-08-01 for the July issue.
%% This is the AUTHORED July issue kept as the monthly exemplar: replace the prose,
%% keep the section order (signal / what changed / analytics / emerging-fading /
%% trial watch / gaps + direction / pooled register + provenance).
%% Build: author claude/build/monthly.tex, then run_all.py --rollup monthly --date <period-end>.
\documentclass[10pt]{article}
\providecommand{\FIGDIR}{fig}
\newcommand{\DIGESTDATE}{31 July 2026}
% ===================== PM Research Watch - shared preamble =====================
\usepackage[T1]{fontenc}
\usepackage[utf8]{inputenc}
\usepackage{amsmath,amssymb}
\usepackage{tgpagella}
\usepackage{tgheros}
\renewcommand{\sfdefault}{qhv}
\renewcommand{\ttdefault}{lmtt}
\usepackage{microtype}
\usepackage[english]{babel}

\usepackage[a4paper,top=17mm,bottom=17mm,left=17mm,right=17mm,
            headsep=6mm,footskip=10mm]{geometry}
\usepackage{graphicx}
\usepackage{xcolor}
\usepackage{colortbl}
\usepackage{tikz}
\usetikzlibrary{calc}
\usepackage[most]{tcolorbox}
\usepackage{titlesec}
\usepackage{fancyhdr}
\usepackage{booktabs}
\usepackage{array}
\usepackage{longtable}
\usepackage{enumitem}
\usepackage{multicol}
\usepackage{ragged2e}
\usepackage[hidelinks]{hyperref}

% ------------------------------------------------------------------ palette
\definecolor{Ink}    {HTML}{14262E}
\definecolor{Deep}   {HTML}{0F4C5C}
\definecolor{Teal}   {HTML}{1F7A8C}
\definecolor{Sky}    {HTML}{4FA3C4}
\definecolor{Sage}   {HTML}{6FA287}
\definecolor{Amber}  {HTML}{E3A008}
\definecolor{Clay}   {HTML}{B5643C}
\definecolor{Coral}  {HTML}{D1495B}
\definecolor{Violet} {HTML}{6A5B9E}
\definecolor{Slate}  {HTML}{4A6572}
\definecolor{Paper}  {HTML}{FBF9F5}
\definecolor{Mist}   {HTML}{EEEAE2}
\definecolor{Rule}   {HTML}{D9D3C9}

\pagecolor{Paper}
\color{Ink}

\hypersetup{colorlinks=true, linkcolor=Deep, urlcolor=Teal, citecolor=Teal,
            pdftitle={PM Research Watch}, pdfauthor={Automated daily digest}}

% ------------------------------------------------------------------ headings
\titleformat{\section}
  {\sffamily\bfseries\large\color{Deep}}{}{0pt}
  {}[{\vspace{1.2mm}\color{Amber}\titlerule[1.1pt]}]
\titlespacing*{\section}{0pt}{5.5mm}{2.8mm}

\titleformat{\subsection}
  {\sffamily\bfseries\normalsize\color{Teal}}{}{0pt}{}
\titlespacing*{\subsection}{0pt}{3.6mm}{1.4mm}

\setlist[itemize]{leftmargin=4.2mm, itemsep=1.1mm, topsep=1.3mm, parsep=0pt,
                  label=\textcolor{Teal}{\footnotesize$\blacktriangleright$}}

% ------------------------------------------------------------------ header/footer
\pagestyle{fancy}
\fancyhf{}
\renewcommand{\headrulewidth}{0pt}
\renewcommand{\footrulewidth}{0pt}
\fancyhead[L]{\sffamily\scriptsize\color{Slate}PM RESEARCH WATCH\;\textcolor{Rule}{\textbar}\;\DIGESTDATE}
\fancyhead[R]{\sffamily\scriptsize\color{Slate}Particulate matter monitoring \& health effects}
\fancyfoot[C]{\sffamily\scriptsize\color{Slate}\thepage}

% ------------------------------------------------------------------ components
% metric tile
\newcommand{\metric}[3]{%
  \begin{tcolorbox}[enhanced, colback=white, colframe=Rule, boxrule=0.5pt,
    arc=1.6mm, left=2.2mm, right=2.2mm, top=1.8mm, bottom=1.8mm,
    borderline west={1.7mm}{0pt}{#3}, nobeforeafter, valign=center,
    height=17mm, width=\linewidth]
    \RaggedRight
    {\sffamily\bfseries\fontsize{17}{17}\selectfont\color{#3}#1}\\[0.5mm]
    {\sffamily\fontsize{6.5}{7.4}\selectfont\color{Slate}#2}
  \end{tcolorbox}}

% pull-quote / key insight box
\newtcolorbox{keybox}[1][]{enhanced, colback=Mist, colframe=Mist,
  arc=1.6mm, left=3.2mm, right=3.2mm, top=2.4mm, bottom=2.4mm,
  borderline west={1.7mm}{0pt}{Amber}, #1}

% subtopic band
\newcommand{\band}[2]{%
  \vspace{3.2mm}\noindent
  \begin{tcolorbox}[enhanced, colback=#2, colframe=#2, arc=1.2mm,
    left=2.6mm, right=2.6mm, top=1.1mm, bottom=1.1mm, boxrule=0pt,
    width=\linewidth, nobeforeafter]
    {\sffamily\bfseries\small\color{white}#1}
  \end{tcolorbox}\vspace{1.6mm}}

% one paper entry: {short cite}{title}{journal}{doi}{body}
\newcommand{\paperentry}[5]{%
  \par\vspace{1.9mm}\noindent
  {\sffamily\bfseries\footnotesize\color{Ink}#1}\;%
  {\sffamily\fontsize{7}{8}\selectfont\color{Slate}#3}\\[0.4mm]
  {\footnotesize\itshape\color{Deep}#2}\\[0.6mm]
  {\footnotesize #4}\\[0.3mm]
  {\fontsize{6.6}{7.6}\selectfont\color{Slate}\texttt{doi:}\,\href{https://doi.org/#5}{\texttt{\detokenize{#5}}}}
  \par}

% caption style
\newcommand{\figcap}[1]{\par\vspace{1.0mm}
  {\fontsize{7.2}{8.6}\selectfont\color{Slate}\textsf{\textbf{Fig.}}\; #1\par}\vspace{1.5mm}}

\providecommand{\permil}{\textperthousand}

\setlength{\parindent}{0pt}
\setlength{\parskip}{0pt}
\renewcommand{\arraystretch}{1.18}


\begin{document}

% ==================================================================== MASTHEAD
\thispagestyle{empty}
\begin{tikzpicture}[remember picture, overlay]
  \fill[Deep] (current page.north west) rectangle
      ($(current page.north east)+(0,-67mm)$);
  \fill[Coral] ($(current page.north west)+(0,-67mm)$) rectangle
      ($(current page.north east)+(0,-68.4mm)$);
  \foreach \x/\y/\r in {22/12/0.9, 41/26/1.5, 63/9/0.7, 88/31/1.1, 112/17/1.9,
                        138/29/0.8, 159/13/1.3, 178/54/1.0, 196/21/1.6,
                        30/44/0.6, 74/57/1.2, 125/48/0.9, 168/6/0.8, 191/59/0.7,
                        52/18/1.0, 100/60/0.6, 148/20/1.1, 66/49/1.4, 205/45/0.8}
    \fill[Sky, opacity=0.30] ($(current page.north west)+(\x mm,-\y mm)$) circle (\r mm);
\end{tikzpicture}

\vspace*{4mm}
{\sffamily\fontsize{7.6}{9}\selectfont\color{Coral}%
 MONTHLY ROLLUP\;\textbullet\;PhD RESEARCH WATCH\par}
\vspace{2.2mm}
{\sffamily\bfseries\fontsize{27}{29}\selectfont\color{white}%
 Particulate Matter\par}
{\sffamily\fontsize{27}{29}\selectfont\color{Sky}%
 Monitoring \& Health Effects\par}
\vspace{5.0mm}
{\sffamily\fontsize{8.4}{10}\selectfont\color{white}%
 July 2026\;\;\textcolor{Coral}{\textbar}\;\;Entry window 1--31 July 2026\;\;%
 \textcolor{Coral}{\textbar}\;\;146 records pooled\par}
\vspace{18mm}

% -------------------------------------------------------------- metric strip
\noindent
\begin{minipage}[t]{0.190\linewidth}\metric{146}{RECORDS POOLED\\(10 SUBTOPICS)}{Deep}\end{minipage}\hfill
\begin{minipage}[t]{0.190\linewidth}\metric{54}{EFFECT ESTIMATES\\EXTRACTED}{Teal}\end{minipage}\hfill
\begin{minipage}[t]{0.190\linewidth}\metric{23}{SENSING /\\INSTRUMENTATION}{Sky}\end{minipage}\hfill
\begin{minipage}[t]{0.190\linewidth}\metric{44}{RECORDS FROM CHINA\\(30\% OF CORPUS)}{Amber}\end{minipage}\hfill
\begin{minipage}[t]{0.190\linewidth}\metric{8}{ULTRAFINE /\\NANOSCALE}{Clay}\end{minipage}

\vspace{6mm}

% ------------------------------------------------------------ signal of month
\section*{\sffamily\bfseries\color{Deep}Signal of the month}
\vspace{-1.5mm}{\color{Coral}\rule{\linewidth}{1.1pt}}\vspace{2.5mm}

\begin{keybox}
{\sffamily\bfseries\small\color{Deep}July was the month the field stopped treating
PM$_{2.5}$ mass as the exposure of record. Three independent literatures ---
ultrafine number, chemical constituents, and unmodelled condensable mass --- each
put a number on what bulk mass misses.}\\[1.4mm]
\footnotesize
Lelieveld et al.\ (\textit{Cardiovasc Res}) built a 1\,km global ultrafine-particle field
from Earth observation and machine learning, found urban annual means of
\textbf{10{,}000--30{,}000 particles\,cm$^{-3}$}, and paired it with a pooled cohort
hazard ratio to estimate \textbf{1.99 (0.81--3.89) million excess deaths per year} ---
roughly 5\% of non-communicable-disease mortality, about half cardiovascular, with
\textbf{91\% occurring in urban and suburban environments}. They put an annual limit of
5{,}000\,cm$^{-3}$ at a \textbf{45\% reduction} in that excess. UFPs are regulated
almost nowhere and measured by no mass monitor. Su et al.\ (\textit{Environ Sci Technol})
supplied the second number: industrial \emph{condensable} PM, absent from every
operational source-apportionment model, contributes a median \textbf{3.6\% (95\% CI
0.9--14.5\%)} of urban PM$_{2.5}$ in four Chinese cities and \textbf{9.5\% in winter},
and its omission inflates apparent secondary inorganic aerosol by \textbf{5.6--16.6\%}.
Easton et al.\ (\textit{Environ Int}) supplied the mechanism: cruise-ship ultrafine PM,
enriched in vanadium, nickel and cobalt, was markedly more pro-inflammatory than
same-port PM of other origin \emph{and} suppressed the innate antiviral response,
increasing epithelial permissiveness to rhinovirus-16 and SARS-CoV-2 --- an effect
reproduced by vanadium alone. Read together with Zhu et al.\ (29 Jul issue), where
PM$_{2.5}$ was the \emph{null} term in its own six-pollutant time-series, the month's
verdict is unambiguous: \textbf{a mass-only monitor is now measuring the wrong quantity
for a growing fraction of the health signal.}
\end{keybox}

\vspace{3mm}

\section*{\sffamily\bfseries\color{Deep}What changed in July}
\vspace{-1.5mm}{\color{Coral}\rule{\linewidth}{1.1pt}}\vspace{2.5mm}

\footnotesize
\begin{itemize}
\item \textbf{Sensing was the single largest subtopic of the month (23/146, 16\%), and
that inverts the picture the daily issues gave.} The 30 July and 1 August issues each
recorded \emph{zero} sensing records and flagged it as a suspected harvest defect. The
month proves the defect: instrumentation work was being published throughout July and
the narrow daily windows simply missed it. No conclusion about the sensing literature
should be drawn from any single daily issue in this series.

\item \textbf{Low-cost sensors were validated in July with unusual honesty about where
they fail.} Gharibzadeh et al.\ needed \emph{seasonal} multivariable calibration
(humidity, temperature, PM$_1$) to reach R$^2$ 0.62--0.93 in Tehran, and reported
persistent underestimation before it. Ali et al.\ found accuracy collapsing with particle
size --- \textbf{R$^2$ $\approx$ 0.90--0.95 for PM$_1$ but 0.28--0.68 for PM$_{10}$} ---
and, critically, that neither the PMS7003 nor the SPS30 captured the dominant
sub-0.30\,$\mu$m fraction at all. Coefficient of variation exceeded 43\% in the coarse
bins. The message for network design is that proprietary size-bin conversions should not
be trusted above 1\,$\mu$m.

\item \textbf{Sensor-plus-satellite fusion produced the month's most striking
performance number, and it deserves scepticism.} Jung et al.\ reported a 10-fold CV
\textbf{R$^2$ of 0.98 (RMSE 1.23\,$\mu$g/m$^3$)} for a 250\,m daily PM$_{2.5}$ field over
Taiwan built from 16{,}795 low-cost sensors plus MAIAC AOD. That is far above the
0.82 (Singh et al., North-East India) and 0.41--0.80 (Yu et al., Shenzhen sidewalks)
obtained elsewhere in the same month. With 16{,}795 sensors as both predictor and
validation target, random 10-fold splits leak spatially; the rolling-origin and spatial
folds are the numbers to read, not the headline.

\item \textbf{Exposure misclassification moved from caveat to quantified bias.} Zhang
et al.\ (\textit{Environ Int}) showed nearest-monitor assignment is \emph{systematically
higher} than both scaled-matched estimates and true personal exposure from outdoor
sources, for every pollutant tested. Nyoni et al.\ (1 Aug issue) showed road-geometry
proxies inverting in sign between African countries. Kos Sinhala et al.\ showed in-home
black carbon driven by open fireplaces and wood stoves rather than ambient concentration.
Three independent designs, one conclusion: the exposure variable in most published
epidemiology is not the exposure that occurred.

\item \textbf{Reproductive and developmental evidence consolidated around
vulnerability, not average effect.} Sieber et al.\ found PM$_{2.5}$ associated with
\textbf{6\% (95\% CI 0--11\%) higher odds of term low birth weight per 10\,$\mu$g/m$^3$}
across 1.4 million Spanish births, but maternal education dominated it --- primary versus
tertiary education carried \textbf{OR 1.76 (1.71--1.82)}. Meeting the WHO 5\,$\mu$g/m$^3$
guideline would have prevented 2{,}032 (95\% CI $-$74 to 3{,}900) cases over nine years;
universal tertiary education would have prevented far more. That comparison is the
uncomfortable and important finding.

\item \textbf{Two well-powered nulls are worth as much as the positives.} Sun et al.\
found \textbf{no} adverse association between PM$_{2.5}$ and FEV$_1$ decline in US adults
with cystic fibrosis (+0.002 z-score per 3\,$\mu$g/m$^3$; 95\% CI $-$0.001, 0.005) at a
mean exposure of 9.1\,$\mu$g/m$^3$. Tien-Smith et al.\ found no clear evidence for
PM$_{2.5}$ and multiple sclerosis (1.01; 0.77--1.32) or motor neuron disease (1.07;
0.86--1.33), while confirming Parkinson's disease risk in Asia (1.19; 1.01--1.41).
Neither is a negative result about air pollution; both are results about where the
exposure-response function is flat.
\end{itemize}

\vspace{3mm}

% ==================================================================== FIGURES
\section*{\sffamily\bfseries\color{Deep}Corpus analytics --- July 2026}
\vspace{-1.5mm}{\color{Coral}\rule{\linewidth}{1.1pt}}\vspace{2.5mm}

\noindent\includegraphics[width=\linewidth]{\FIGDIR/f1_subtopics.png}
\figcap{Subtopic distribution across all 146 pooled records. All ten subtopics are
populated --- no daily issue in July achieved that. Sensing, forecasting \&
instrumentation leads at 23 records, which is the direct refutation of the zero counts
reported in the 30 July and 1 August daily issues. Assignment is single-label by
dominant contribution.}

\vspace{3mm}

\noindent\includegraphics[width=\linewidth]{\FIGDIR/w1_subtopic_trend.png}
\figcap{Records per issue date by subtopic. \textbf{Read this as a harvest diagnostic,
not as a time series.} The 26 July bar is the retrospective backfill for 1--26 July
harvested in a single batch on 1 August; it is not a day's output. Only the 27--30 July
bars are genuine daily windows. The panel is included because the contrast is the
point: four narrow daily windows recovered 95 records while one month-wide query
recovered 51 more that the dailies never saw.}

\clearpage

\noindent\includegraphics[width=\linewidth]{\FIGDIR/f2_design_endpoint.png}
\figcap{Left --- study architecture over the month. Measurement campaigns are the
largest class (30), ahead of observational cohorts (25), modelling/inventory (22) and
reviews (22). Only one trial or intervention record appeared in the entire month, which
is the standing weakness of this evidence base: almost nothing here is experimental in
humans. Right --- health endpoint distribution.}

\vspace{3mm}

\noindent\includegraphics[width=\linewidth]{\FIGDIR/f3_metric_geography.png}
\figcap{Left --- particle metric. PM$_{2.5}$ alone in 75 records and PM$_{2.5}$ with
PM$_{10}$ in 32; \textbf{ultrafine or nanoscale metrics appear in only 8 records (5\%)}
despite carrying the month's largest single mortality estimate. That mismatch between
where the burden is claimed and where the measurement effort goes is the clearest
instrumentation gap in this corpus. Right --- geography. China supplies 44 records
(30\%), Europe 22, North America 20; sub-Saharan Africa supplies 4 and Latin America 5.}

\clearpage

\noindent\includegraphics[width=\linewidth]{\FIGDIR/f4_heatmap.png}
\figcap{Subtopic $\times$ study-architecture cross-tabulation over the month. The dense
Sensing/measurement-campaign and Burden/modelling cells are expected. The informative
absence is the near-empty trial column across every health subtopic.}

\vspace{3mm}

\noindent\includegraphics[width=\linewidth]{\FIGDIR/f5_forest.png}
\figcap{All 54 quantitative effect estimates extracted across July on a common ratio
axis. \textbf{This is not a meta-analysis and no pooled summary is drawn.} Exposure
contrasts are not commensurable --- per-IQR, per-5, per-10\,$\mu$g/m$^3$, quartile and
quintile contrasts, and mixture indices all appear --- and the outcome scale mixes RR,
OR, aOR and HR across different diseases. \textbf{Heterogeneity is therefore not
estimable and none is reported}; the one place a pooled figure would be legitimate is
Tien-Smith et al.'s within-paper neurodegenerative meta-analysis, which is reproduced
here as three separate points rather than absorbed. The visible structure is that
per-10\,$\mu$g/m$^3$ PM$_{2.5}$ estimates cluster tightly just above unity (1.02--1.16)
while the large ratios belong to per-IQR contrasts, mixture indices and gaseous
co-pollutants.}

\clearpage

\noindent\includegraphics[width=\linewidth]{\FIGDIR/f6_lifecourse.png}
\figcap{Life-course windows addressed across the month. Counts are non-exclusive and
cover only records with an identifiable human life-course window; \textit{in vitro},
ex-vivo, animal and pure-instrumentation records are excluded. Adolescence remains the
thinnest window in the corpus, as it was in every individual daily issue.}

\vspace{4mm}

\section*{\sffamily\bfseries\color{Deep}Emerging, fading, and recurring}
\vspace{-1.5mm}{\color{Coral}\rule{\linewidth}{1.1pt}}\vspace{2.5mm}

\footnotesize
\textbf{Emerging.} Three themes appeared or firmed up materially this month.
\textbf{(i) Ultrafine particles as a regulatory object} --- Lelieveld et al.\ and a
paired \textit{Cardiovasc Res} commentary, Easton et al.'s ship-source toxicology,
Williams et al.'s hepatic xanthine-oxidoreductase mechanism, Zeng et al.'s microglial
circRNA work, and Olivares Lopez et al.'s finding that particle-number simulations are
strongly sensitive to emission-inventory choice all converge in a single month, driven by
the new EU Ambient Air Quality Directive's particle-number monitoring requirement.
\textbf{(ii) Compound exposures} --- PM$_{2.5}$ with heat is now a design rather than a
sensitivity analysis (Yang et al., 610{,}613 OHCAs; Zhang et al., 5.07 million adults).
\textbf{(iii) Omics as the mediation layer} --- proteomic signatures for type 2 diabetes
(49{,}134 UK Biobank participants), gut-microbiome mediation of cognition, cord-blood
extracellular-vesicle miRNA, and placental mtDNA copy number and telomere length.

\vspace{2mm}
\textbf{Fading.} Nothing can honestly be called fading on one month of data, and the
backfill/daily imbalance described in the provenance box would manufacture false
declines if pressed. The single defensible observation is that \textbf{single-pollutant,
single-endpoint time-series studies now read as the low-information tail} of this corpus:
several appeared, most were single-city and single-hospital, and the ones that reported a
PM term alongside gases generally found PM the weakest term in its own model.

\vspace{2mm}
\textbf{Recurring groups and venues.} No author cluster recurs meaningfully at
first-author level --- the apparent repeats (Li 5, Zhang 4, Wang 4, Chen 4) are surname
frequency, not the same investigators, and should not be reported as a cluster. At
institutional level, Chinese public-health schools (Zhengzhou, Sichuan, Fudan, CUHK) and
the Barcelona/BiSC and Danish/Nordic cohort infrastructures each contribute repeatedly.
Venue concentration is modest: \textit{Toxics} (7), \textit{J Environ Sci (China)} (6),
\textit{Environ Pollut} (5), \textit{Ecotoxicol Environ Saf} (5), \textit{Environ Sci
Technol} (4), \textit{J Hazard Mater} (4). The month's highest-standing placements were
\textit{Eur Heart J}, \textit{J Am Coll Cardiol}, \textit{Cardiovasc Res},
\textit{Environ Health Perspect} (3 records), \textit{Environ Int} (3) and
\textit{Int J Epidemiol} (2). Citation counts on records days to weeks old are
uninformative and were not used to rank anything here.

\clearpage

% ------------------------------------------------------------------ trial watch
\section*{\sffamily\bfseries\color{Deep}Trial registry watch}
\vspace{-1.5mm}{\color{Coral}\rule{\linewidth}{1.1pt}}\vspace{2.5mm}

\footnotesize
ClinicalTrials.gov returned \textbf{one} PM or air-pollution record registered or
updated inside the monitored portion of the window: \textbf{NCT05160948}, ``Assessing
Changes in Knowledge, Behavior, and Attitude Regarding Air Pollution After Receiving
Personal Air Monitoring Results'' (Cincinnati Children's Hospital Medical Center;
observational; $n = 105$; completed). No endpoint analysis was run because the study is
observational rather than interventional. It is nonetheless the single most
on-topic registry item this watch has seen: it tests whether \emph{returning personal
monitor data to the participant} changes behaviour, which is the assumption every
citizen-science and personal-sensing deployment rests on and almost none evaluates.

\vspace{1.5mm}
Two records in the peer-reviewed corpus function as the interventional evidence the
registry did not supply. Gonz\'alez-Colom et al.\ (\textit{BMJ Open Respir Res},
NCT06421402) monitored household air quality with low-cost sensors in 205 adults with
COPD, asthma or bronchiectasis for two months: \textbf{51.7\% of homes had at least one
pollutant in an at-risk category, and the burden was overwhelmingly PM$_{2.5}$-driven
(40.1\% of dwellings at risk)}, with active or passive smoking strongly associated
($p < 0.001$). \textbf{No relationship was found with hospitalisations or ED visits} ---
the honest null that a two-month screening window cannot overturn. Bolden et al.\
reviewed indoor wildfire-smoke mitigation and concluded that portable-air-cleaner
efficacy is real but too dependent on building envelope and occupant behaviour to
generalise across the national housing stock. \textbf{The interventional evidence base
for personal and household PM reduction remains thin, and July did not thicken it.}

\vspace{4mm}

% -------------------------------------------------------------------- open gaps
\section*{\sffamily\bfseries\color{Deep}Open gaps and read on direction}
\vspace{-1.5mm}{\color{Coral}\rule{\linewidth}{1.1pt}}\vspace{2.5mm}

\begin{keybox}
{\sffamily\bfseries\footnotesize\color{Deep}Gaps in the literature}\\[1.2mm]
\footnotesize
\textbf{(i) Ultrafine measurement does not exist at network scale.} Eight records used an
ultrafine metric; one of them attributes two million annual deaths to it. There is no
low-cost particle-number instrument in this corpus with a published field-calibration
comparable to what Gharibzadeh et al.\ did for PM$_{2.5}$ mass.
\textbf{(ii) Almost no human intervention studies.} One trial-class record in 146.
Every causal claim in this month's corpus rests on observational designs with
modelled exposure.
\textbf{(iii) Sub-0.3\,$\mu$m is invisible to the deployed sensor base.} Ali et al.\
state this directly for the two most widely deployed OPC modules; it is also where the
ultrafine burden lives.
\textbf{(iv) Geographic imbalance is severe.} 44 records from China against 4 from
sub-Saharan Africa, in a month whose most methodologically important exposure paper
(Nyoni et al.) is precisely about African settings where the standard proxies fail.
\end{keybox}

\vspace{2.5mm}

\begin{keybox}
{\sffamily\bfseries\footnotesize\color{Deep}Gaps this pipeline creates}\\[1.2mm]
\footnotesize
\textbf{(i) No date-windowed channel into \textit{Atmos.\ Meas.\ Tech.} or
\textit{Atmos.\ Chem.\ Phys.}} PubMed does not index them, the Consensus connector has
no entry-date filter, and OpenAlex date filtering now returns 403/429. Every
instrumentation record in this issue reached the corpus through a
biomedically-indexed journal, so the pure atmospheric-measurement literature is
\emph{still} absent even at 23 sensing records. The fix remains a Crossref by-journal
feed keyed on \texttt{created} for AMT (1867-8548) and ACP (1680-7324).
\textbf{(ii) Narrow daily windows systematically under-recover.} Demonstrated
quantitatively this month: the sensing subtopic read 0, 0, 1 and 6 across four
consecutive daily issues, against 23 for the month.
\textbf{(iii) The 1--26 July backfill is classification-depth only.} Effect estimates
were extracted for that batch only where the abstract reported an interval; a genuine
daily cadence would have captured more.
\end{keybox}

\vspace{3mm}

\footnotesize
\textbf{Read on direction.} The field is moving from \emph{mass} to \emph{composition and
number}, and the instrumentation base has not moved with it. July produced a
two-million-death annual estimate for a pollutant metric that no regulatory network
measures and no validated low-cost sensor reports; simultaneously it produced careful
evidence that the deployed optical sensor base cannot see below 0.3\,$\mu$m and degrades
badly above 1\,$\mu$m. That gap is the most fundable problem visible in this corpus, and
it is an instrumentation problem before it is an epidemiological one. \textbf{The
falsifiable prediction: within the next two monthly issues, expect at least one
field-calibration study of a low-cost particle-number or black-carbon sensor against a
CPC or aethalometer reference, framed explicitly against the new EU particle-number
monitoring requirement.} If that does not appear by the September issue, the
measurement community is lagging the epidemiology by more than the usual cycle, and the
right inference is that the bottleneck is hardware cost rather than research attention.

\clearpage

% ==================================================================== TABLE
\section*{\sffamily\bfseries\color{Deep}Pooled corpus register --- July 2026}
\vspace{-1.5mm}{\color{Coral}\rule{\linewidth}{1.1pt}}\vspace{2.5mm}

{\fontsize{6.7}{8.4}\selectfont
\setlength{\tabcolsep}{2.4pt}
\begin{longtable}{@{}>{\RaggedRight}p{25mm}>{\RaggedRight}p{41mm}>{\RaggedRight}p{27mm}%
>{\RaggedRight}p{23mm}>{\RaggedRight}p{18mm}>{\RaggedRight}p{15mm}@{}}
\toprule
\textsf{\textbf{First author}} & \textsf{\textbf{Focus}} & \textsf{\textbf{Journal}} &
\textsf{\textbf{Design}} & \textsf{\textbf{Metric}} & \textsf{\textbf{Region}}\\
\midrule
\endfirsthead
\toprule
\textsf{\textbf{First author}} & \textsf{\textbf{Focus}} & \textsf{\textbf{Journal}} &
\textsf{\textbf{Design}} & \textsf{\textbf{Metric}} & \textsf{\textbf{Region}}\\
\midrule
\endhead
\multicolumn{6}{@{}l@{}}{\cellcolor{Mist}\textsf{\textbf{\textcolor{Deep}{Neuro / mental health}}}}\\[0.6mm]
\href{https://doi.org/10.1001/jamanetworkopen.2026.25015}{Dales} & Wildfire-sourced PM$_{2.5}$ and ED visits for migraine and other primary headache & \textit{JAMA Netw Open} & Case-crossover & PM$_{2.5}$ (wildfire) & Canada\\
\href{https://doi.org/10.1093/ageing/afag215}{Glover} & Umbrella meta-review of urban environmental exposures and cognitive health & \textit{Age Ageing} & Umbrella review / meta-analysis & PM$_{2.5}$ & Global\\
\href{https://doi.org/10.1016/j.nsa.2026.107019}{Goossens} & In utero air pollution exposure and pre-/postnatal brain development & \textit{Neurosci Appl} & Narrative review & PM$_{2.5}$ / PM$_{10}$ / PM$_{\mathrm{coarse}}$ & Global\\
\href{https://doi.org/10.1016/j.isci.2026.116617}{Ma} & PM$_{2.5}$-induced neurotoxicity and depression: exposure risks and mechanisms & \textit{iScience} & Narrative review & PM$_{2.5}$ + components & Global\\
\href{https://doi.org/10.3760/cma.j.cn112150-20250804-00752}{Pan} & Maternal PM exposure and offspring neurodevelopment: effects and mechanisms & \textit{Zhonghua Yu Fang Yi Xue Za Zhi} & Narrative review & PM (bulk) & China\\
\addlinespace[0.9mm]
\multicolumn{6}{@{}l@{}}{\cellcolor{Mist}\textsf{\textbf{\textcolor{Teal}{Cardiovascular \& metabolic}}}}\\[0.6mm]
\href{https://doi.org/10.3760/cma.j.cn112338-20251120-00843}{Huang} & PM$_{2.5}$ and dyslipidaemia in rural Chinese children; dietary effect modification & \textit{Zhonghua Liu Xing Bing Xue Za Zhi} & Cross-sectional (multistage) & PM$_{2.5}$ & China\\
\href{https://doi.org/10.32394/pe/218966}{Kazmierska} & Air pollution and cardiovascular disease: review of current data & \textit{Przegl Epidemiol} & Narrative review & PM$_{2.5}$ / PM$_{10}$ & Global\\
\href{https://doi.org/10.1007/s00125-026-06814-2}{Xi} & Air pollution, built environment and temperature vs. prediabetes phenotypes & \textit{Diabetologia} & Prospective cohort & PM$_{2.5}$ & Germany\\
\href{https://doi.org/10.1161/JAHA.125.046917}{Yang} & Short-term PM$_{2.5}$ and hospitalisation/ED use in atrial fibrillation & \textit{J Am Heart Assoc} & Case-crossover & PM$_{2.5}$ (wildfire-influenced) & USA\\
\addlinespace[0.9mm]
\multicolumn{6}{@{}l@{}}{\cellcolor{Mist}\textsf{\textbf{\textcolor{Sky}{Reproductive \& developmental}}}}\\[0.6mm]
\href{https://doi.org/10.1016/j.jes.2025.10.043}{Chen} & Prenatal air pollution, fetal growth trajectory phenotypes and child neurodevelopment & \textit{J Environ Sci (China)} & Prospective cohort & PM$_{2.5}$ / PM$_{10}$ & China\\
\href{https://doi.org/10.1016/j.ecoenv.2026.120566}{Huang} & Joint prenatal bisphenol + air-pollutant mixture exposure and preschool IQ & \textit{Ecotoxicol Environ Saf} & Prospective cohort (mixtures) & PM$_{2.5}$ / PM$_{10}$ & China\\
\href{https://doi.org/10.1007/s40572-026-00552-8}{Peterson} & Ambient air pollution, PAHs and pubertal development: critical review & \textit{Curr Environ Health Rep} & Critical review & PM$_{2.5}$ components & Global\\
\href{https://doi.org/10.1016/j.envres.2026.125313}{Xie} & Long-term PM$_{2.5}$ and timing of pubertal development in boys and girls & \textit{Environ Res} & Prospective cohort & PM$_{2.5}$ & China\\
\addlinespace[0.9mm]
\multicolumn{6}{@{}l@{}}{\cellcolor{Mist}\textsf{\textbf{\textcolor{Sage}{Respiratory \& allergic}}}}\\[0.6mm]
\href{https://doi.org/10.1007/s00484-026-03289-z}{Bayar Muluk} & Climate change and otolaryngology: airway disease impacts and emission reduction & \textit{Int J Biometeorol} & Narrative review & PM$_{2.5}$ & Global\\
\href{https://doi.org/10.4103/mgr.MEDGASRES-D-25-00164}{Liao} & Short-term air pollution, respiratory symptoms and lung function in asthma & \textit{Med Gas Res} & Panel study & PM$_{2.5}$ / PM$_{10}$ & China\\
\href{https://doi.org/10.3390/epidemiologia7040097}{Sanchez de Prada} & Ambient PM and asthma-coded admissions in children 0-3 y, Spain & \textit{Epidemiologia} & Ecological panel & PM$_{2.5}$ / PM$_{10}$ & Spain\\
\href{https://doi.org/10.1016/j.jhazmat.2026.143076}{Yang} & Allergenic vs. toxic risk of airborne seed hairs compared with ambient PM$_{2.5}$ & \textit{J Hazard Mater} & Field measurement & PM$_{2.5}$ + bioaerosol & China\\
\addlinespace[0.9mm]
\multicolumn{6}{@{}l@{}}{\cellcolor{Mist}\textsf{\textbf{\textcolor{Violet}{Mechanistic toxicology}}}}\\[0.6mm]
\href{https://doi.org/10.1016/j.jes.2025.12.074}{Ding} & PM$_{2.5}$ promotes vascular calcification via miR-27b-5p/HES1 signalling & \textit{J Environ Sci (China)} & In vitro + organoid + murine & PM$_{2.5}$ & China\\
\href{https://doi.org/10.1007/s00360-026-01698-5}{Maraschi} & Osmoregulatory disruption and metal bioaccumulation from settleable atmospheric PM & \textit{J Comp Physiol B} & Controlled exposure (ecotox) & Settleable APM & Brazil\\
\href{https://doi.org/10.1016/j.jes.2025.11.018}{Wu} & Macrophage NLRP3 inflammasome in real-ambient PM-induced hepatic glucose dysregulation & \textit{J Environ Sci (China)} & Real-ambient murine + co-culture & Real-ambient PM & China\\
\addlinespace[0.9mm]
\multicolumn{6}{@{}l@{}}{\cellcolor{Mist}\textsf{\textbf{\textcolor{Clay}{Occupational \& indoor}}}}\\[0.6mm]
\href{https://doi.org/10.3389/fmedt.2026.1796828}{Chen} & Nested suction instrument to cut intraoperative surgical-smoke PM$_{2.5}$ & \textit{Front Med Technol} & Bench / device experiment & PM$_{2.5}$ (surgical smoke) & China\\
\href{https://doi.org/10.1093/annweh/wxag055}{Cole} & Determinants of respirable crystalline silica exposure in Australian tunnelling & \textit{Ann Work Expo Health} & Observational exposure study & Respirable dust / RCS & Australia\\
\href{https://doi.org/10.1016/j.jhazmat.2026.143081}{Ding} & Stage-resolved hair Zn-isotope dynamics as a firefighting combustion-exposure tracer & \textit{J Hazard Mater} & Repeated-measures biomarker & Combustion-derived PM & China\\
\href{https://doi.org/10.3390/jox16040119}{Ramsden} & Xenobiotic hazards in aircraft cabin air & \textit{J Xenobiot} & Review + risk economics & Ultrafine particles & Global\\
\addlinespace[0.9mm]
\multicolumn{6}{@{}l@{}}{\cellcolor{Mist}\textsf{\textbf{\textcolor{Amber}{Sensing, forecasting \& instrumentation}}}}\\[0.6mm]
\href{https://doi.org/10.1016/j.jenvman.2026.130515}{Shi} & Adaptive spatiotemporal graph transformer for multi-site multi-step PM$_{2.5}$ forecasting & \textit{J Environ Manage} & Deep-learning model & PM$_{2.5}$ & China / India\\
\href{https://doi.org/10.3390/bios16070380}{Tsuchiya} & Seasonal and indoor-outdoor characteristics of fluorescent aerosol particles in offices & \textit{Biosensors} & Real-time sensor campaign & Bioaerosol / coarse & Japan\\
\addlinespace[0.9mm]
\multicolumn{6}{@{}l@{}}{\cellcolor{Mist}\textsf{\textbf{\textcolor{Amber}{Exposure assessment \& modelling}}}}\\[0.6mm]
\href{https://doi.org/10.1029/2025GH001621}{Solorzano-Araque} & Multiscale exposure modelling: LOTOS-EUROS + mobility surveys + particle toxicity & \textit{GeoHealth} & CTM + mobility integration & PM$_{2.5}$ / PM$_{10}$ & Colombia\\
\addlinespace[0.9mm]
\multicolumn{6}{@{}l@{}}{\cellcolor{Mist}\textsf{\textbf{\textcolor{Coral}{Burden, policy \& mitigation}}}}\\[0.6mm]
\href{https://doi.org/10.1002/hsr2.72886}{Karamzad} & Ischaemic heart disease burden in MENA 1990-2021; PM attributable fraction 28.2\% & \textit{Health Sci Rep} & GBD secondary analysis & Ambient PM & MENA\\
\href{https://doi.org/10.1371/journal.pone.0347989}{Rees} & HEPA filtration to reduce respiratory infection in care homes & \textit{PLoS One} & RCT process evaluation & Filtered PM & UK\\
\href{https://doi.org/10.1016/j.jes.2026.01.009}{Shan} & Policy-driven emission reductions and post-pandemic rebound, high-resolution inventories & \textit{J Environ Sci (China)} & Emission inventory & PM emissions & China\\
\addlinespace[0.9mm]
\multicolumn{6}{@{}l@{}}{\cellcolor{Mist}\textsf{\textbf{\textcolor{Slate}{Other clinical endpoints}}}}\\[0.6mm]
\href{https://doi.org/10.3390/dj14070440}{Lupita} & PM$_{2.5}$, preventive-care capacity and caries/edentulism burden across the EU & \textit{Dent J} & Ecological panel & PM$_{2.5}$ & EU-27\\
\href{https://doi.org/10.1016/j.jdin.2026.06.008}{Park} & Long-term PM exposure and risk of androgenetic alopecia & \textit{JAAD Int} & Nationwide cohort & PM$_{2.5}$ / PM$_{10}$ & South Korea\\
\href{https://doi.org/10.3171/2026.3.PEDS2615}{Sun} & Environmental exposure and area deprivation vs. paediatric hydrocephalus outcomes & \textit{J Neurosurg Pediatr} & Registry cohort + geospatial & PM$_{2.5}$ (satellite-derived) & USA\\
\addlinespace[0.9mm]
\bottomrule
\end{longtable}}


\vspace{4mm}
\begin{keybox}
{\sffamily\bfseries\footnotesize\color{Deep}Provenance and method --- read this before
citing any count above}\\[1.2mm]
{\fontsize{7.2}{8.8}\selectfont
\textbf{This issue breaks the standing rollup contract, deliberately and on request.}
Rollups in this series normally aggregate previously published daily issues and never
re-query an API. Daily coverage of July exists only for 27--30 July, so the 1--26 July
portion of this corpus was \textbf{harvested independently on 1 August 2026} and had
never appeared in any daily issue. Sources for that backfill: \textbf{PubMed}
E-utilities on the \texttt{[Date - Entry]} field, 2026/07/01--2026/07/26, along two
independent axes --- health (352 records) and instrumentation (53) --- giving 368 unique
records after merge, of which 22 were already in \texttt{state/seen.json} from the daily
issues and were dropped, 166 passed a PM-relevance title screen, and \textbf{51 were
selected for full classification} on the basis of journal standing, methodological
weight and subtopic coverage. Those 51 are written to
\texttt{state/corpus/2026-07-26.json} with \texttt{backfill: true} and an
\texttt{entry\_window} recording the retrospective harvest.
\\[1.0mm]
\textbf{The consequence for every count in this issue, stated plainly:} the month's 146
records are \textbf{not} a complete census of July. They are 95 records from four
exhaustively-screened daily windows plus 51 selected from 166 eligible records across
the remaining 26 days. \textbf{The 1--26 July portion is therefore a curated sample, not
a census, and subtopic proportions for that period are biased toward what was selected.}
Specifically, sensing and exposure work was preferentially retained because it is the
reader's primary field. The 23-record sensing count is a \emph{floor}, not an estimate,
and the correct reading of the headline finding is that sensing output in July was
\emph{at least} 23 records and certainly more --- which strengthens rather than weakens
the conclusion that the daily windows were under-recovering. No proportion in this issue
should be quoted as a July prevalence.
\\[1.0mm]
The daily-trend panel is labelled a harvest diagnostic for the same reason: the 26 July
bar is one batch, not one day. Records were deduplicated against
\texttt{state/seen.json} on PMID, lowercased DOI and title signature before any
summarisation, so no record appears twice across the daily issues and the backfill.
Effect estimates were transcribed verbatim from abstracts and only where an interval was
reported; 54 qualified. Subtopic, design, particle-metric and geography labels are
single-label judgements from title, abstract, keywords and MeSH terms, not a controlled
vocabulary. Two records fetched without abstracts (\textit{Am J Epidemiol} PMID 42429684;
\textit{J Am Coll Cardiol} PMID 42455106) were \textbf{excluded rather than summarised
from their titles}, despite the JACC piece being directly on this issue's theme. DOIs
carry through from PubMed's own record. Content derived from PubMed, Europe PMC and
ClinicalTrials.gov metadata; abstracts remain the copyright of their respective
publishers.\par}
\end{keybox}

\end{document}
